\documentclass[tikz,border=6pt]{standalone}
\usepackage{peabra}
\begin{document}
\begin{tikzpicture}[x=1cm,y=1cm,
  lep/.style={draw=peKek, fill=peHalvany, rounded corners=2pt, line width=0.8pt, align=center,
    font=\footnotesize, text width=30mm, minimum height=13mm, inner sep=3pt},
  szam/.style={circle, fill=peKiemel, text=white, font=\scriptsize\bfseries, inner sep=0, minimum size=5mm}]

  \node[lep] (l1) at (0.00,0.00) {A rendszer és határainak rögzítése};
  \node[szam] at (l1.north west) {1};
  \node[lep] (l2) at (3.90,0.00) {A V lépéseinek felírása a saját rendszerre};
  \node[szam] at (l2.north west) {2};
  \node[lep] (l3) at (7.80,0.00) {Munkatermék minden lépéshez};
  \node[szam] at (l3.north west) {3};
  \node[lep] (l4) at (11.70,0.00) {Jobb ági pár és ellenőrzési mód minden bal ági lépéshez};
  \node[szam] at (l4.north west) {4};
  \node[lep] (l5) at (0.00,-2.20) {Felülvizsgálati pontok belépési feltételekkel};
  \node[szam] at (l5.north west) {5};
  \node[lep] (l6) at (3.90,-2.20) {Felelősök és határidők};
  \node[szam] at (l6.north west) {6};
  \node[lep] (l7) at (7.80,-2.20) {A terv karbantartása a projekt során};
  \node[szam] at (l7.north west) {7};
  \draw[nyil, line width=0.7pt] (l1.east) -- (l2.west);
  \draw[nyil, line width=0.7pt] (l2.east) -- (l3.west);
  \draw[nyil, line width=0.7pt] (l3.east) -- (l4.west);
  \draw[nyil, line width=0.7pt] (l4.south) -- ++(0,-0.38) -| (l5.north);
  \draw[nyil, line width=0.7pt] (l5.east) -- (l6.west);
  \draw[nyil, line width=0.7pt] (l6.east) -- (l7.west);
  \draw[-{Stealth[length=2mm]}, draw=pePiros, line width=0.7pt, dash pattern=on 3pt off 2pt] ($(l7.north)+(-0.55,0)$) to[bend left=18] node[midway, left=1pt, font=\scriptsize, text=pePiros, align=right] {új tudás, változás} ($(l3.south)+(-0.55,0)$);
\end{tikzpicture}
\end{document}
